%-------------------------
% Resume in Latex
% Author : Sourabh Bajaj
% Website: https://github.com/sb2nov/resume
% License : MIT
%------------------------

\documentclass[letterpaper,11pt]{article}

\usepackage{latexsym}
\usepackage[empty]{fullpage}
\usepackage{titlesec}
\usepackage{marvosym}
\usepackage[usenames,dvipsnames]{color}
\usepackage{verbatim}
\usepackage{enumitem}
\usepackage[pdftex]{hyperref}
\usepackage{fancyhdr}


\pagestyle{fancy}
\fancyhf{} % clear all header and footer fields
\fancyfoot{}
\renewcommand{\headrulewidth}{0pt}
\renewcommand{\footrulewidth}{0pt}

% Adjust margins
\addtolength{\oddsidemargin}{-0.375in}
\addtolength{\evensidemargin}{-0.375in}
\addtolength{\textwidth}{1in}
\addtolength{\topmargin}{-.5in}
\addtolength{\textheight}{1.0in}

\urlstyle{same}

\raggedbottom
\raggedright
\setlength{\tabcolsep}{0in}

% Sections formatting
\titleformat{\section}{
  \vspace{-4pt}\scshape\raggedright\large
}{}{0em}{}[\color{black}\titlerule \vspace{-5pt}]

%-------------------------
% Custom commands
\newcommand{\resumeItem}[2]{
  \item\small{
    \textbf{#1}{: #2 \vspace{-2pt}}
  }
}

\newcommand{\resumeSubheading}[4]{
  \vspace{-1pt}\item
    \begin{tabular*}{0.97\textwidth}{l@{\extracolsep{\fill}}r}
      \textbf{#1} & #2 \\
      \textit{\small#3} & \textit{\small #4} \\
    \end{tabular*}\vspace{-5pt}
}

\newcommand{\resumeSubItem}[2]{\resumeItem{#1}{#2}\vspace{-4pt}}

\renewcommand{\labelitemii}{$\circ$}

\newcommand{\resumeSubHeadingListStart}{\begin{itemize}[leftmargin=*]}
\newcommand{\resumeSubHeadingListEnd}{\end{itemize}}
\newcommand{\resumeItemListStart}{\begin{itemize}}
\newcommand{\resumeItemListEnd}{\end{itemize}\vspace{-5pt}}

%-------------------------------------------
%%%%%%  CV STARTS HERE  %%%%%%%%%%%%%%%%%%%%%%%%%%%%


\begin{document}

%----------HEADING-----------------
\begin{tabular*}{\textwidth}{l@{\extracolsep{\fill}}r}
  \textbf{\href{https://www.linkedin.com/in/edgar-ivan-hinojosa-salda\%C3\%B1a-95273b260/}{\Large Edgar Ivan Hinojosa Saldaña}} & Email : \href{e_e.ivan.h_hs@outlook.com}{e\_e.ivan.h\_hs@outlook.com}\\
	\href{ https://edgarivanhinojosa.xyz}{ \small https://edgarivanhinojosa.xyz}\
  \href{http://sourabhbajaj.com/}{} & Mobile : +52 867 175 128 8 \\
\end{tabular*}



Profesional con iniciativa propia y aprendizaje autónomo, enfocado en la precisión y la mejora continua de procesos.
En mi desarrollo he adquirido una serie de técnicas y habilidades en modelado matemático, análisis estadístico de datos, desarrollo de software e instrumentación.
%-----------EDUCATION-----------------
\section{Educación}
  \resumeSubHeadingListStart
    \resumeSubheading
      {Instituto America de Estudios Superiores}{Nuevo Laredo, Tamps.}
      {Preparatoria }{Ago. 2018 -- Ago. 2021}

    \resumeSubheading
      {Facultad de Ciencias Físico Matemáticas, UANL}{San Nicolás de los Garza }
      {Licenciatura en Física }{Ago. 2021 -- Dic. 2025}
		\resumeItemListStart
        \resumeItem{Cursos Relevantes}
          {Metodología de Programación, Física Computacional, Probabilidad y Estadística, Modelado Científico }
      \resumeItemListEnd
  \resumeSubHeadingListEnd
%-----------Distinciones-----------------
\section{Distinciones}
  \resumeSubHeadingListStart
    \resumeSubheading
  {Olimpiada Mexicana de Matemáticas}{OMM}
	  {Finalista estatal. Resolución de problemas y razonamiento matemático}{ 2019-2021}

    \resumeSubheading
      {Programa de Talentos}{UANL}
      {Programa de becados por excelencia académica}{2021}
 \resumeItemListEnd
  \resumeSubHeadingListEnd


%-----------I+D-----------------
\section{Investigación y Desarrollo}
  \resumeSubHeadingListStart
    \resumeSubheading
  {XXIV Verano de Investigación Científica y Tecnológica }{UANL}
	  {Documentación de demostraciones deductivas de óptica geométrica}{ Ago 2023}

\resumeSubheading
	{Asistente de Investigación en Simulación de Materiales}{CICFIM}
	{Desarrollo de algoritmos de optimización en Computador de Alto Rendimiento}{2022-2023}

\resumeSubheading
	{Asistente de Investigación en Físico Matemática}{CICFIM}
	{Deducción con motivación física de estructura simplectica en Mecánica Clásica}{2023-2024}

    \resumeSubheading
  {XXVI Verano de Investigación Científica y Tecnológica }{UANL}
	  {Diseño de máquina de vórtices para sistemas hidrodinámicos}{ Ago 2025}

    \resumeSubheading
  {Escritura de Tesina}{UANL}
	  {Documentación de cálculos, revisión bibliográfica, interpretación de resultados}{ Ago 2023}


 \resumeItemListEnd
  \resumeSubHeadingListEnd





%-----------PROJECTS-----------------
\section{Modelado Científico y Análisis de Datos}
 \textit{Destrezas y habilidades desarrolladas}

\resumeItemListStart
\resumeItem{Desarrollo y validación de modelos matemáticos}{Deducidos a primeros principios o esquemas estadísticos; contrastados experimentalmente para su interpretación y sucesiva corrección.}
\resumeItem{Revisión de Literatura}{Lectura de documentación y antecedentes bibliográficos, así como escritura de entragables expositivos, reportería interactiva.}
\resumeItem{Desarrollo Software}{Implementación de rutinas en Python para limpieza, análisis y visualización de datos experimentales. Algoritmos de extracción de patrones con visión computacional semi automática}

\resumeItem{Técnicas estadísticas}{Intervalos de confianza, pruebas de hipótesis, pruebas de normalidad, regresión lineal y logística, validación cruzada, bootstrapping, arboles de decisión, bosques aleatorios, dbscan, clustering, maquinas de soporte vectorial, metodos de regularización, lasso, análisis de componentes principales, redes neuronales entre otros.}
\resumeItemListEnd

 \textit{Proyectos}

\resumeItemListStart
   \resumeSubItem{Distribución de la aprobación de Profesores}
      {Aplicación de pruebas de normalidad, Kolmogorov, entre otras para carterización de distribición de aprobación de profesorado por los usuarios de misprofes.com}
    \resumeSubItem{Librería de Métodos Numéricos con Fortran Moderno }
      {Una librería de Fortran y Python usando Fortran Moderno con Orientación a objetos para la organización de métodos númericos y estadísticos.}
   \resumeSubItem{Extracción de haces láser}
      {Procesamiento de imagenes para la extracción de haces laser por método de optimización de serpientes, motivado físicamente.}
   \resumeSubItem{Codigo Predictivo en Python y Fortran}
      {Implementación computacional de una red neuronal con verosimilitud biológica con respuesta a estímulos estocásticos, para aprendizaje de distribuciones de probabilidad. }
   \resumeSubItem{Modelado fraccional de la Viscolásticidad de una Gelatina}
      {Validación experimental de modelo por medio de análisis de resultados experimentales de respuestas dinámicas}
   \resumeSubItem{Mantenimiento Predictivo}
      {Arboles de decisión y Bosques aleatorios para caracterización de vibración en motores de potencia para su sucesiva predicción }

\resumeItemListEnd


  \resumeSubHeadingListEnd
%
%----------COURSES&CERTIFICATIONS------
\section{Cursos y Certificaciones}
  \resumeSubHeadingListStart
\resumeSubheading
      {Programación en Python: Nivel Básico e Intermedio}{UANL}
	  {Fundamentos de Lenguaje de Programación en Python}{ Ene 2022}

    \resumeSubheading
{Escuela de Verano: Modelación y Herramientas para el Análisis de Datos, ed. II }{UANL}
	  {Bash Scripting, Density Functional Theory, Dinámica Molecular}{ Jul 2022}
\resumeSubheading
      {Fundamentos de AutoCAD Autodesk}{UANL}
	  {Introducción al uso de software para dibujo tecnico}{ Oct 2022}

 \resumeSubheading
      {Escuela de Verano: Modelación y Herramientas para el Análisis de Datos, ed. III}{FCFM}
	  {Lenguaje de programación R, Métodos Estadísticos, Bash Scripting}{Jul. 2023}

\resumeSubheading
      {Taller de Consultoria}{FCFM}
	  {Métodos estadísticas aplicados, reportería para ejecutivos}{Jun. 2024}


\resumeSubheading
      {Sexta Escuela Temática en Análisis y Sistemas Dinámicos}{CIMAT, Gto}
	  {Teoría espectral en análisis armónico, sistemas dinámicos caóticos}{Jun. 2024}

\resumeSubheading
      {Escuela de Verano en Computo Estadístico 2024}{CIMAT, Mty}
	  {Aprendizaje maquina, servicios de consultoría, computo paralelo }{Jul. 2024}

\resumeSubheading{Taller de Seguridad Láser}{CICESE, MTY} 
{Introducción a normatividad aplicable y protocolos de seguridad en laboratorio. }

\resumeSubheading
{Escuela en Tecnologías Avanzadas e Integradas 2025}{CICESE, MTY}
{Tecnologías ópticas, teoría de plasmones }{2025}


\resumeSubheading
{Statistical Learning}{StanfordOnline, MOOC}
 {Modelos lineales y de clasificación, redes neuronales convolucionales, bootstraping y validación cruzada.}

\resumeSubheading
      {Escuela de Verano: Modelación y Herramientas para el Análisis de Datos, ed. V}{UANL}
	  {Introducción a FeynCalc para cálculo simbólico en teoría cuántica de campos}{ Jul 2023}

  \resumeSubHeadingListEnd

%-----------Congresos y coloquios-----------------
\section{Congresos y coloquios}

  \resumeSubHeadingListStart
    \resumeSubItem{Congreso Nacional de Física 2024}
	{Asistencia y exposición del poster: \textit{Geometría Simpléctica en Mecánica Clásica}}

    \resumeSubItem{Coloquio de Asociación Americana de Profesores de Física 2024}
	{Ponencia de charla: \textit{Experimento de Eddington un precedente experimental hidrodinámico}}

    \resumeSubItem{Seminario de Óptica Aplicada y Sustentabilidad  de Energía 2025}
	{Asistencia y Exposición de poster (Segundo Lugar): \textit{Medición de Fuerza de Fricción con Análisis de Series de Tiempo} }

    \resumeSubItem{Feria de Física Experimental FCFM Primavera 2025}
	{Exposición de poster: \textit{La Lente Vorticial de un Fluido y su Similitud con la Lente Gravitacional }}

    \resumeSubItem{Congreso Nacional de Física 2025}
	{Exposición de poster: \textit{Analogía de lente gravitacional por guía de
onda en sistema de fluidos con vorticidad} }



  \resumeSubHeadingListEnd


  %--------ACTIVIDADES EXTRACURRICULARES----------------------
\section{Actividades extracurriculares}
  \resumeSubHeadingListStart
    \resumeSubItem{Datathon 2024 TEC-MTY}
      {Web Scrapping, Sentimental Anlysis, Clustering}
    \resumeSubItem{Jurado de Feria de Experimentos 2025}
      {Participación como Juez en Feria de Experimentas con enfoque de experimento didáctico y propuesta tecnologica}


  \resumeSubHeadingListEnd
%



%-----------EXPERIENCE-----------------
\section{Experiencia}
  \resumeSubHeadingListStart


\resumeSubheading
{Creativeit Group}{Mar. 2026 -- Presente}
{Ingeniero de Pruebas de Validación }{(medio tiempo)}
\resumeItemListStart
\resumeItem{Diseño de sistema de adquisición de datos}{ programada con Python y PyVisa, usando los protocolos SCIP de la IEE 484, e integración de interfaz gráfica con PyQT.}
\resumeItem{Documentación y ejecución de pruebas de validación}{Consulta de estándares y escritura de informes para pre certificación de productos y ensayos para mantenimiento preventivo. }
\resumeItemListEnd

\resumeSubheading
{Departamento de Física, FCFM}{Ene. 2025 -- Dic. 2025}
{Técnico y Docente de Laboratorio}{}
\resumeItemListStart
\resumeItem{Diseño e instrucción de prácticas de laboratorio}{ enfatizando el análisis cuantitativo y estadístico de resultados.}
\resumeItem{Diseño y captura de base de datos}{ con SQLite para inventariado de equipos almacenados. }
\resumeItemListEnd

  \resumeSubHeadingListEnd
%--------SKILLS----------------------
\section{Habilidades}
 \resumeSubHeadingListStart
   \item{
     \textbf{Idiomas}{: Español, Inglés  }
   }
\item {
     \textbf{Lenguajes de Programación}{: R, Python, Bash, Fortran, Mathematica} \ \hfill
}

   \item{
\textbf{Herramientas:} Jupyter, interfaces gráficas (streamlit, PyQT), procesamiento de imágenes (opencv), Numpy, Scikit-Learn, PyTorch
}
\item {
     \textbf{Software}{: Office, AutoCAD, RStudio, Onshape, Vim, Entornos Linux}
}
\item{
\textbf{Instrumentación:} Osciloscopio, multímetro de precisión, electrómetro, probadores de alta tensión, analizador vectorial de redes, sistemas de adquisición de datos, láseres, sistemas ópticos,  fuentes de alto voltaje, soldadura, circuitos y electrónica digital. \\

 \resumeSubHeadingListEnd


%-------------------------------------------
\end{document}

